% Algebra 2 Lesson 6-1 fast-track packet: Key Features of Exponential Functions
% (enVision A2 2024 TE, a2_6-1_TE.tex; Example 2 and Example 5 only, chosen for Topic 6 test items 1, 4, 17 and Lesson Quiz items 1, 5)
% Re-cut 2026-10-31 toward the Topic 6 test (was Examples 1, 3).
% Shared preamble for the "fast" (hacked-down) Algebra 2 lesson packets, Oct 2026.
% Each packet: % Shared preamble for the "fast" (hacked-down) Algebra 2 lesson packets, Oct 2026.
% Each packet: % Shared preamble for the "fast" (hacked-down) Algebra 2 lesson packets, Oct 2026.
% Each packet: \input{fast-packet-template} then \begin{document} ... \end{document}.
\documentclass[12pt]{article}
\usepackage[margin=0.7in]{geometry}
\usepackage{amsmath,amssymb}
\usepackage{booktabs,tabularx,array}
\usepackage{enumitem}
\usepackage[T1]{fontenc}
\usepackage[utf8]{inputenc}
\usepackage{lmodern}
\usepackage{tikz}
\usepackage{pgfplots}
\pgfplotsset{compat=1.18}
\usepackage{xcolor}
\setlist{nosep,leftmargin=1.4em}
\setlength{\parindent}{0pt}
\setlength{\parskip}{6pt}
\renewcommand{\arraystretch}{1.35}
\newcolumntype{Y}{>{\raggedright\arraybackslash}X}
\pagestyle{empty}

% blanks and workspace
\newcommand{\blank}[1][2.6cm]{\rule[-2pt]{#1}{0.5pt}}
\newcommand{\lines}[1][1]{\par\vspace{2pt}\foreach \i in {1,...,#1}{\rule{0pt}{22pt}\rule{\linewidth}{0.4pt}\par}}
\newcommand{\wspart}[2]{\vspace{4pt}{\large\bfseries #1: #2}\par\vspace{-6pt}\rule{\textwidth}{0.4pt}\par}

% header: \packethead{Lesson 1-3}{Piecewise-Defined Functions}
\newcommand{\packethead}[2]{%
{\LARGE\bfseries #2}\hfill Name \blank[5cm]\\[2pt]
Algebra 2 \quad #1 (fast track) \hfill Date \blank[2.5cm] \quad Period \blank[1.2cm]\par\vspace{2pt}}

% boxed goals: \goals{learning goal}{language goal}
\newcommand{\goals}[2]{%
\fbox{\parbox{\dimexpr\textwidth-2\fboxsep-2\fboxrule}{%
\textbf{Learning goal:} #1\\[2pt]
\textbf{Language goal:} #2}}\par}

% quiz target box: \quiztarget{text}
\newcommand{\quiztarget}[1]{%
\vspace{4pt}\colorbox{gray!12}{\parbox{\dimexpr\textwidth-2\fboxsep}{\textbf{On the quiz:} #1}}\par}

% IXL table row helper: code & skill & which quiz question it feeds
\newcommand{\ixlrow}[3]{\texttt{#1} & #2 & #3 \\}

% one grid style for every graph; override xmin/xmax/... per plot
\pgfplotsset{sheetgrid/.style={
  axis lines=middle, grid=both, grid style={gray!45},
  xmin=-10.5, xmax=10.5, ymin=-10.5, ymax=10.5,
  xtick={-10,-8,...,10}, ytick={-10,-8,...,10},
  x=0.28cm, y=0.28cm,
  ticklabel style={font=\scriptsize, fill=white, inner sep=1pt},
  xlabel={$x$}, ylabel={$y$},
  axis line style={thick}, clip=true}}
% unit grid for small windows
\pgfplotsset{unitgrid/.style={sheetgrid,
  xmin=-5.5, xmax=5.5, ymin=-5.5, ymax=5.5,
  xtick={-5,...,5}, ytick={-5,...,5}, x=0.5cm, y=0.5cm}}
% open and closed endpoints
\tikzset{closed/.style={circle, fill, inner sep=1.6pt},
         open/.style={circle, draw, fill=white, inner sep=1.4pt}}
 then \begin{document} ... \end{document}.
\documentclass[12pt]{article}
\usepackage[margin=0.7in]{geometry}
\usepackage{amsmath,amssymb}
\usepackage{booktabs,tabularx,array}
\usepackage{enumitem}
\usepackage[T1]{fontenc}
\usepackage[utf8]{inputenc}
\usepackage{lmodern}
\usepackage{tikz}
\usepackage{pgfplots}
\pgfplotsset{compat=1.18}
\usepackage{xcolor}
\setlist{nosep,leftmargin=1.4em}
\setlength{\parindent}{0pt}
\setlength{\parskip}{6pt}
\renewcommand{\arraystretch}{1.35}
\newcolumntype{Y}{>{\raggedright\arraybackslash}X}
\pagestyle{empty}

% blanks and workspace
\newcommand{\blank}[1][2.6cm]{\rule[-2pt]{#1}{0.5pt}}
\newcommand{\lines}[1][1]{\par\vspace{2pt}\foreach \i in {1,...,#1}{\rule{0pt}{22pt}\rule{\linewidth}{0.4pt}\par}}
\newcommand{\wspart}[2]{\vspace{4pt}{\large\bfseries #1: #2}\par\vspace{-6pt}\rule{\textwidth}{0.4pt}\par}

% header: \packethead{Lesson 1-3}{Piecewise-Defined Functions}
\newcommand{\packethead}[2]{%
{\LARGE\bfseries #2}\hfill Name \blank[5cm]\\[2pt]
Algebra 2 \quad #1 (fast track) \hfill Date \blank[2.5cm] \quad Period \blank[1.2cm]\par\vspace{2pt}}

% boxed goals: \goals{learning goal}{language goal}
\newcommand{\goals}[2]{%
\fbox{\parbox{\dimexpr\textwidth-2\fboxsep-2\fboxrule}{%
\textbf{Learning goal:} #1\\[2pt]
\textbf{Language goal:} #2}}\par}

% quiz target box: \quiztarget{text}
\newcommand{\quiztarget}[1]{%
\vspace{4pt}\colorbox{gray!12}{\parbox{\dimexpr\textwidth-2\fboxsep}{\textbf{On the quiz:} #1}}\par}

% IXL table row helper: code & skill & which quiz question it feeds
\newcommand{\ixlrow}[3]{\texttt{#1} & #2 & #3 \\}

% one grid style for every graph; override xmin/xmax/... per plot
\pgfplotsset{sheetgrid/.style={
  axis lines=middle, grid=both, grid style={gray!45},
  xmin=-10.5, xmax=10.5, ymin=-10.5, ymax=10.5,
  xtick={-10,-8,...,10}, ytick={-10,-8,...,10},
  x=0.28cm, y=0.28cm,
  ticklabel style={font=\scriptsize, fill=white, inner sep=1pt},
  xlabel={$x$}, ylabel={$y$},
  axis line style={thick}, clip=true}}
% unit grid for small windows
\pgfplotsset{unitgrid/.style={sheetgrid,
  xmin=-5.5, xmax=5.5, ymin=-5.5, ymax=5.5,
  xtick={-5,...,5}, ytick={-5,...,5}, x=0.5cm, y=0.5cm}}
% open and closed endpoints
\tikzset{closed/.style={circle, fill, inner sep=1.6pt},
         open/.style={circle, draw, fill=white, inner sep=1.4pt}}
 then \begin{document} ... \end{document}.
\documentclass[12pt]{article}
\usepackage[margin=0.7in]{geometry}
\usepackage{amsmath,amssymb}
\usepackage{booktabs,tabularx,array}
\usepackage{enumitem}
\usepackage[T1]{fontenc}
\usepackage[utf8]{inputenc}
\usepackage{lmodern}
\usepackage{tikz}
\usepackage{pgfplots}
\pgfplotsset{compat=1.18}
\usepackage{xcolor}
\setlist{nosep,leftmargin=1.4em}
\setlength{\parindent}{0pt}
\setlength{\parskip}{6pt}
\renewcommand{\arraystretch}{1.35}
\newcolumntype{Y}{>{\raggedright\arraybackslash}X}
\pagestyle{empty}

% blanks and workspace
\newcommand{\blank}[1][2.6cm]{\rule[-2pt]{#1}{0.5pt}}
\newcommand{\lines}[1][1]{\par\vspace{2pt}\foreach \i in {1,...,#1}{\rule{0pt}{22pt}\rule{\linewidth}{0.4pt}\par}}
\newcommand{\wspart}[2]{\vspace{4pt}{\large\bfseries #1: #2}\par\vspace{-6pt}\rule{\textwidth}{0.4pt}\par}

% header: \packethead{Lesson 1-3}{Piecewise-Defined Functions}
\newcommand{\packethead}[2]{%
{\LARGE\bfseries #2}\hfill Name \blank[5cm]\\[2pt]
Algebra 2 \quad #1 (fast track) \hfill Date \blank[2.5cm] \quad Period \blank[1.2cm]\par\vspace{2pt}}

% boxed goals: \goals{learning goal}{language goal}
\newcommand{\goals}[2]{%
\fbox{\parbox{\dimexpr\textwidth-2\fboxsep-2\fboxrule}{%
\textbf{Learning goal:} #1\\[2pt]
\textbf{Language goal:} #2}}\par}

% quiz target box: \quiztarget{text}
\newcommand{\quiztarget}[1]{%
\vspace{4pt}\colorbox{gray!12}{\parbox{\dimexpr\textwidth-2\fboxsep}{\textbf{On the quiz:} #1}}\par}

% IXL table row helper: code & skill & which quiz question it feeds
\newcommand{\ixlrow}[3]{\texttt{#1} & #2 & #3 \\}

% one grid style for every graph; override xmin/xmax/... per plot
\pgfplotsset{sheetgrid/.style={
  axis lines=middle, grid=both, grid style={gray!45},
  xmin=-10.5, xmax=10.5, ymin=-10.5, ymax=10.5,
  xtick={-10,-8,...,10}, ytick={-10,-8,...,10},
  x=0.28cm, y=0.28cm,
  ticklabel style={font=\scriptsize, fill=white, inner sep=1pt},
  xlabel={$x$}, ylabel={$y$},
  axis line style={thick}, clip=true}}
% unit grid for small windows
\pgfplotsset{unitgrid/.style={sheetgrid,
  xmin=-5.5, xmax=5.5, ymin=-5.5, ymax=5.5,
  xtick={-5,...,5}, ytick={-5,...,5}, x=0.5cm, y=0.5cm}}
% open and closed endpoints
\tikzset{closed/.style={circle, fill, inner sep=1.6pt},
         open/.style={circle, draw, fill=white, inner sep=1.4pt}}

\begin{document}

\packethead{Lesson 6-1}{Key Features of Exponential Functions}
\goals{I can tell how a reflection or a shift moves the $y$-intercept and the asymptote of an exponential function, and compare two exponential functions by their growth or decay factors.}{I can explain an exponential function using the words \emph{growth factor}, \emph{decay factor}, \emph{y-intercept}, \emph{asymptote}, and \emph{reflection}.}

\quiztarget{(1) Example 2: which function is $3^x$ after a vertical stretch by 8 and a reflection across the $x$-axis? (2) Growth or decay after 5 years (full Savvas quiz only). (3, bonus) Domain, range, $y$-intercept, asymptote of $\tfrac12\cdot6^x$ (Example 2's rules). (4) Write and graph a decay model (full Savvas quiz only). (5) Example 5: which function has a greater growth factor than a graph?}
\topictarget{6}{Question 1 asks which function is $5^x$ stretched by 2 and reflected across the $x$-axis (Example 2). Question 4 asks which function decays faster than a graphed curve (Example 5). Question 17 asks for the $y$-intercept and asymptote of $-4(0.3)^x+1$ (Example 2; Form B's $2(3)^{-x}$ is Example 1, which this packet skips).}

{\footnotesize Words: exponential function (\emph{funci\'on exponencial}), growth factor (\emph{factor de incremento}), decay factor (\emph{factor de decremento}). An \textbf{exponential function} is $f(x)=a\cdot b^x$ with $a\neq0$, $b>0$, $b\neq1$; its $y$-intercept is $a$ because $b^0=1$.\par}

\wspart{Example 2}{Graph Transformations of Exponential Functions}
Graph each function. Describe the graph in terms of transformations of the parent function $f(x)=3^x$. How do the asymptote and $y$-intercept of the given function compare to the asymptote and $y$-intercept of the parent function?

\begin{minipage}[t]{0.49\textwidth}
\vspace{0pt}
\textbf{A.} $g(x)=-3^x$

\begin{center}
\begin{tikzpicture}
% TE window x -5 to 5, y -5 to 5, labels -4,-2,2,4 on both axes, plus half a unit so the arrowheads show
\begin{axis}[unitgrid, xticklabels={,-4,,-2,,,,2,,4,}, yticklabels={,-4,,-2,,,,2,,4,},
  ylabel style={xshift=-19pt}, x=0.35cm, y=0.35cm]
\addplot[very thick, blue, domain=-5.5:1.55, samples=80, <->] {3^x};
\addplot[very thick, red, domain=-5.5:1.55, samples=80, <->] {-(3^x)};
\node[closed, blue] at (axis cs:0,1) {};
\node[closed, red] at (axis cs:0,-1) {};
\node[font=\scriptsize, inner sep=1pt, anchor=north east] at (axis cs:0,0) {$O$};
\node[font=\scriptsize, fill=white, inner sep=1pt, anchor=west] at (axis cs:0.25,0.75) {$(0,1)$};
\node[font=\scriptsize, fill=white, inner sep=1pt, anchor=west] at (axis cs:0.25,-0.75) {$(0,-1)$};
\node[font=\scriptsize, blue, fill=white, inner sep=1pt] at (axis cs:2.4,4.6) {$f$};
\node[font=\scriptsize, red, fill=white, inner sep=1pt] at (axis cs:2.4,-4.6) {$g$};
\end{axis}
\end{tikzpicture}
\end{center}

$g(x)=-3^x=-f(x)$: multiplying by $-1$ \textbf{reflects} the graph across the $x$-axis.
\begin{itemize}[topsep=0pt]\small
\item $y$-intercept: moves from $a$ to $-a$, so from $1$ to $-1$
\item Asymptote: does not change; still the $x$-axis ($y=0$)
\end{itemize}
\end{minipage}\hfill
\begin{minipage}[t]{0.49\textwidth}
\vspace{0pt}
\textbf{B.} $h(x)=3^x-4$

\begin{center}
\begin{tikzpicture}
% TE window x -5 to 5, y -5 to 5, labels -4,-2,2,4; f approaches y=0 and h approaches y=-4 on the left
\begin{axis}[unitgrid, xticklabels={,-4,,-2,,,,2,,4,}, yticklabels={,-4,,-2,,,,2,,4,},
  ylabel style={xshift=-19pt}, x=0.35cm, y=0.35cm]
\addplot[very thick, blue, domain=-5.5:1.55, samples=80, <->] {3^x};
\addplot[very thick, red, domain=-5.5:2.04, samples=80, <->] {3^x-4};
\node[closed, blue] at (axis cs:0,1) {};
\node[closed, red] at (axis cs:0,-3) {};
\node[font=\scriptsize, inner sep=1pt, anchor=north east] at (axis cs:0,0) {$O$};
\node[font=\scriptsize, fill=white, inner sep=1pt, anchor=south east] at (axis cs:-0.5,0.7) {$(0,1)$};
\node[font=\scriptsize, fill=white, inner sep=1pt, anchor=west] at (axis cs:0.4,-3) {$(0,-3)$};
\node[font=\scriptsize, blue, fill=white, inner sep=1pt] at (axis cs:0.6,4.6) {$f$};
\node[font=\scriptsize, red, fill=white, inner sep=1pt] at (axis cs:3.0,4.6) {$h$};
\end{axis}
\end{tikzpicture}
\end{center}

$h(x)=3^x-4=f(x)-4$: adding a constant $k$ \textbf{shifts} the graph vertically by $k$ units (here down 4).
\begin{itemize}[topsep=0pt]\small
\item $y$-intercept: moves from $1$ to $1+k$, so to $-3$
\item Asymptote: also moves; it is $y=k$, or $y=-4$
\end{itemize}
\end{minipage}

\textbf{Say it:} For $h(x)=3^x-4$ the \emph{y-intercept} moves from $1$ to \blank[1.2cm] and the \emph{asymptote} moves from $y=0$ to \blank[1.6cm]. For $g(x)=-3^x$ the \emph{reflection} moves the \emph{y-intercept} to \blank[1.2cm], but the \emph{asymptote} \blank[3cm].

\newpage

\wspart{Try It 2}{Compare the asymptote and $y$-intercept}
How do the asymptote and $y$-intercept of the given function compare to the asymptote and $y$-intercept of the function $f(x)=5^x$?

Hint: $f(x)=5^x$ has $y$-intercept $1$ and asymptote $y=0$. Put $x=0$ to find each new $y$-intercept.

\textbf{a.} $g(x)=5^{x+3}$

$y$-intercept: $g(0)=\blank[3cm]$ \qquad Asymptote: $y=\blank[2cm]$

Compared with $f$: \lines[1]

\textbf{b.} $h(x)=5^{-x}+6$

$y$-intercept: $h(0)=\blank[3cm]$ \qquad Asymptote: $y=\blank[2cm]$

Compared with $f$ (which transformations?): \lines[2]

\wspart{Example 5}{Compare Two Exponential Functions}
A museum purchased a painting and a sculpture in the same year. Their values are changing exponentially as shown. Which art work's value is increasing more quickly?

\begin{minipage}[t]{0.47\textwidth}
\vspace{0pt}
\textbf{Sculpture value:} $f(x)=50(1.075)^x$, in thousands of dollars in $x$ years.

\vspace{4pt}
\textbf{Painting value:}

\begin{tikzpicture}
% TE painting graph: Years 0 to 5 by 1, Value (thousands of dollars) 0 to 80 by 20; points (0,40), (1,44), (2,48.4), (5,64.4); no arrowheads
\begin{axis}[sheetgrid, xmin=0, xmax=5.3, ymin=0, ymax=84,
  xtick={0,...,5}, ytick={0,20,...,80},
  xlabel={Years}, ylabel={Value (\$1000s)},
  xlabel style={at={(axis description cs:0.5,-0.12)}, anchor=north},
  ylabel style={at={(axis description cs:-0.2,0.5)}, rotate=90, anchor=south},
  axis lines=left, x=1.0cm, y=0.05cm]
\addplot[very thick, red, domain=0:5, samples=60] {40*1.1^x};
\node[closed] at (axis cs:0,40) {};
\node[closed] at (axis cs:1,44) {};
\node[closed] at (axis cs:2,48.4) {};
\node[closed] at (axis cs:5,64.4) {};
\node[font=\scriptsize, red, fill=white, inner sep=1pt, anchor=north west] at (axis cs:0.2,35) {$(0,40)$};
\node[font=\scriptsize, red, fill=white, inner sep=1pt, anchor=north west] at (axis cs:1.08,42) {$(1,44)$};
\node[font=\scriptsize, red, fill=white, inner sep=1pt, anchor=south east] at (axis cs:1.95,51) {$(2,48.4)$};
\node[font=\scriptsize, red, fill=white, inner sep=1pt, anchor=south east] at (axis cs:4.92,67) {$(5,64.4)$};
\end{axis}
\end{tikzpicture}
\end{minipage}\hfill
\begin{minipage}[t]{0.50\textwidth}
\vspace{0pt}
Identify key features for the sculpture: $f(x)=50(1.075)^x$. The initial purchase price is \$50,000. The \textbf{growth factor} is $1.075$.

\vspace{4pt}
Find the growth factor, $b$, for the painting. The ratio of $y$-values one year apart is $b$:
\[b=\frac{44}{40}=1.1 \quad\text{and}\quad b=\frac{48.4}{44}=1.1\]
\[g(x)=40(1.1)^x\]
Verify for $x=5$: $g(5)=40(1.1)^5\approx40(1.6105)\approx64.4$.

The growth factor for the painting is $1.10$. Since $1.10>1.075$, the painting's value is increasing more rapidly.
\end{minipage}

\textbf{Say it:} The painting's \emph{growth factor} is \blank[1.2cm] because $44\div40=\blank[1.2cm]$; the sculpture's is \blank[1.4cm]. The painting's value grows faster because \blank[3.6cm]. A \emph{decay factor} of $0.6$ makes a value shrink faster than a \emph{decay factor} of $0.75$ because \dots\lines[1]

\newpage

\wspart{Try It 5}{When does the painting pass the sculpture?}
In Example 5, in what year will the value of the painting become greater than the value of the sculpture?

Hint: compare $f(x)=50(1.075)^x$ and $g(x)=40(1.1)^x$ year by year (thousands of dollars, to the nearest tenth).

\begin{center}
\begin{tabular}{|c|c|c|c|}
\hline
$x$ (years) & Sculpture $f(x)$ & Painting $g(x)$ & Which is greater? \\ \hline
$0$ & $50$ & $40$ & sculpture \\
$5$ & & & \\
$8$ & & & \\
$9$ & & & \\
$10$ & & & \\ \hline
\end{tabular}
\end{center}

The value of the painting becomes greater in year \blank[1.6cm], because \lines[2]

\wspart{IXL}{Practice for the quiz}
\textbf{SmartScore 70 or above on each starred ($\bigstar$) skill replaces your quiz score.}

\vspace{2pt}
\begin{tabularx}{\textwidth}{l Y l}
\toprule
\textbf{Code} & \textbf{Skill} & \textbf{Feeds} \\
\midrule
\ixlrow{LWE}{Evaluate exponential functions}{warm-up (Try It 5 table)}
\ixlrow{RSN $\bigstar$}{Transformations of functions}{Quiz 1}
\ixlrow{EFD $\bigstar$}{Domain and range of exponential functions}{Quiz 3 (bonus)}
\ixlrow{2YT $\bigstar$}{Exponential functions over unit intervals}{Quiz 5}
\bottomrule
\end{tabularx}

\wspart{Exit question}{Model with mathematics}
\fbox{\parbox{\dimexpr\textwidth-2\fboxsep-2\fboxrule}{%
A colony of bacteria starts with 50 organisms and quadruples each day. Write a function, $P(t)$, that represents the population of the bacteria after $t$ days. Then find the number of bacteria that will be in the colony after 5 days.

\vspace{4pt}
Starting amount ($y$-intercept): \blank[1.5cm] \quad \emph{Growth factor}: \blank[1.2cm] because \blank[3.6cm]

\vspace{4pt}
$P(t)=\blank[3.5cm]$ \qquad After 5 days: \blank[4cm] bacteria.

\vspace{4pt}
Is this exponential growth or decay? Compare its \emph{growth factor} with the painting's \emph{growth factor} in Example 5.\lines[1]
}}

\end{document}
